% Generated draft: requires mathematical and source review.
\documentclass[11pt]{article}
\usepackage[margin=25mm]{geometry}
\usepackage{fontspec}
\setmainfont{Latin Modern Roman}
\usepackage{amsmath,amssymb,mathtools,unicode-math}
\setmathfont{Latin Modern Math}
\usepackage{enumitem,needspace,xurl,hyperref}
\setlength{\parindent}{0pt}
\setlength{\parskip}{5pt}
\setlength{\emergencystretch}{3em}
\newcommand{\N}{\mathbb{N}}
\newcommand{\Z}{\mathbb{Z}}
\newcommand{\Q}{\mathbb{Q}}
\newcommand{\R}{\mathbb{R}}
\newcommand{\C}{\mathbb{C}}
\newcommand{\F}{\mathbb{F}}
\newcommand{\A}{\mathbb{A}}
\newcommand{\PP}{\mathbb{P}}
\newcommand{\E}{\mathbb{E}}
\newcommand{\eps}{\varepsilon}
\newcommand{\dd}{\,\mathrm{d}}
\begin{document}
\section*{UnsolvedMath ID: 30002015}
\textbf{Source problem code:} catalogue entry
\subsection{Definitions and mathematical statement}
Write \(G^F=\{g\in G:F(g)=g\}\) for the fixed-point group of the endomorphism \(F\) appearing in the intended Deligne--Lusztig construction. Let \(X_P\) denote the parabolic Deligne--Lusztig variety in the source, and put
\[
V_i=H_c^i(X_P,\mathbb Q_\ell),\qquad V=\bigoplus_{i\geq 0}V_i,
\]
where \(H_c^i\) denotes compactly supported \(\ell\)-adic étale cohomology, with its intended \(G^F\)-action. Such cohomology requires a prime \(\ell\) different from the characteristic of the base field. The record does not specify the underlying group, field, endomorphism, parabolic data, or the construction of \(X_P\).

For finite-dimensional representations of a finite group over \(\mathbb Q_\ell\), call two representations \emph{disjoint} if they have no common irreducible constituent. By semisimplicity, this is equivalent to the vanishing of the space of equivariant homomorphisms between them. For a finite group \(W_S\), the notation \(\mathbb Q_\ell W_S\) means its group algebra over \(\mathbb Q_\ell\). The record does not define \(W_S\) or \(S\), so this interpretation remains conditional on identifying the intended group.

\textbf{Requested assertion, subject to unresolved source hypotheses.} For every datum in the intended, but unspecified, class of parabolic Deligne--Lusztig varieties, prove
\[
\operatorname{Hom}_{\mathbb Q_\ell[G^F]}(V_i,V_j)=0
\qquad\text{for all }i,j\geq 0\text{ with }i\ne j,
\]
and prove an isomorphism of \(\mathbb Q_\ell\)-algebras
\[
\operatorname{End}_{\mathbb Q_\ell[G^F]}\!\left(\bigoplus_{i\geq 0}H_c^i(X_P,\mathbb Q_\ell)\right)
\cong \mathbb Q_\ell W_S.
\]
Here the endomorphisms are those of the total representation displayed in the supplied statement. This transcription does not assert the claim for arbitrary parabolics: the admissible data and the meaning of \(W_S\) must first be recovered from the source. Consequently, the supplied record alone does not determine a fully specified conjecture.
\subsection{Short English statement}
The requested claim is that distinct compactly supported cohomology degrees of the intended parabolic Deligne--Lusztig variety share no irreducible \(G^F\)-representation, and that the endomorphism algebra of their direct sum is \(\mathbb Q_\ell W_S\). The construction, admissible hypotheses, and group \(W_S\) are unspecified. Supplied dated curation labels the problem partially solved; this draft does not independently verify that assessment.
\subsection{Sources}
\begin{itemize}
\item Supplied local UnsolvedMath registry record, permanent ID \(30002015\), source problem code \texttt{catalogue entry}, title \emph{Cohomological Disjointness for Parabolic Deligne--Lusztig Varieties}, in the Oberwolfach Reports Open Problems collection. Project registry: \texttt{data/unsolvedmath/problems.json}; supplementary registry: \texttt{catalog/problems.json}.
\item \emph{Representations of Finite Groups} (2012), Oberwolfach Reports, pp.~963--1019, DOI \texttt{10.4171/owr/2012/16}. Supplied report reference: \url{https://doi.org/10.4171/owr/2012/16}. The record attributes the problem to workshop open-problem discussions and supplies no individual proposer or precise location within the report.
\item Supplied literature reference: \emph{Parabolic Deligne--Lusztig varieties}, \url{https://arxiv.org/abs/1110.4863}.
\item Supplied literature reference: \emph{Translation by the full twist}, \url{https://www.sciencedirect.com/science/article/pii/S0001870819305447}.
\item Literature assessment dated exactly \texttt{2026-08-21}: the supplied curation labels the problem partially solved and reports that the endomorphism-algebra conjecture is proved in important cases, including the split type-A disjointness consequence, while the full general pairwise-disjointness and endomorphism statement remains open. This is supplied dated curation derived from the OWR record, not an independently verified literature assessment; its status and cited claims require independent checking.
\end{itemize}
\end{document}
